\section*{Capítulo \ref{ch:g9:statproba} --- Estadística y probabilidad}

\begin{solution}{exo:g9:statproba:1}
\omterm{def:g9:statproba:indicators}{Significar}: $\dfrac{14 + 9 + 17 + 9 + 12 + 11 + 12}{7} = \dfrac{84}{7} = 12$.

\omterm{def:g9:statproba:indicators}{Mediana}: ordenado, $9, 9, 11, 12, 12, 14, 17$; el medio ($4$ th) el valor es
$12$.

\omterm{def:g9:statproba:indicators}{Rango}: $17 - 9 = 8$.
\end{solution}

\begin{solution}{exo:g9:statproba:2}
\omterm{def:g9:statproba:indicators}{Significar}: $\dfrac{19 + 21 + 24 + 18 + 22 + 25 + 19}{7} = \dfrac{148}{7}
\approx 21.1$ grados.

\omterm{def:g9:statproba:indicators}{Mediana}: ordenado, $18, 19, 19, 21, 22, 24, 25$; el $4$ el valor es $21$
grados.
\end{solution}

\begin{solution}{exo:g9:statproba:3}
\omterm{def:g9:statproba:indicators}{Significar}:
\[
\frac{2 \times 0 + 4 \times 1 + 3 \times 2 + 1 \times 3}{10}
= \frac{0 + 4 + 6 + 3}{10} = \frac{13}{10} = 1.3
\text{ goals per match.}
\]
\omterm{def:g9:statproba:indicators}{Mediana}: la serie ordenada es $0,0,1,1,1,1,2,2,2,3$; el $5$ th y $6$ th
los valores son ambos $1$: \omterm{def:g9:statproba:indicators}{mediana} $1$.
\end{solution}

\begin{solution}{exo:g9:statproba:4}
$\P(3) = \frac16$. \quad
$\P(\text{even}) = \frac36 = \frac12$. \quad
$\P(\text{at least } 3) = \P(\{3,4,5,6\}) = \frac46 = \frac23$. \quad
$\P(\text{not } 6) = 1 - \frac16 = \frac56$.
\end{solution}

\begin{solution}{exo:g9:statproba:5}
$\P(\text{white}) = \frac{5}{12}$.

$\P(\text{red or green}) = \frac{4 + 3}{12} = \frac{7}{12}$.

$\P(\text{not white}) = 1 - \frac{5}{12} = \frac{7}{12}$.

Los dos últimos son iguales: ``no blancos'' \emph{es} el evento ``rojo o
verde''.
\end{solution}

\begin{solution}{exo:g9:statproba:6}
El $12$ resultados (\omterm{def:g5:solids:def}{rostro} de la moneda, valor del dado) son igualmente
probable.

cabezas y un $6$: un resultado de $12$: $\frac{1}{12}$.

cabezas o un $6$: los resultados con cabezas ($6$ de ellos) más el resultado
(colas, $6$): $7$ resultados, entonces $\frac{7}{12}$. Contando cada evento y
restar la superposición da lo mismo:
$\frac{6}{12} + \frac{2}{12} - \frac{1}{12} = \frac{7}{12}$.
\end{solution}

\begin{solution}{exo:g9:statproba:7}
Primer sorteo: rojo $\frac25$, negro $\frac35$. Segundo empate sin
reemplazo: después de un rojo, $1$ rojo y $3$ los negros permanecen entre $4$; después
un negro, $2$ rojo y $2$ negro entre $4$.

Dos negros: $\frac35 \times \frac24 = \frac{6}{20} = \frac{3}{10}$.

Diferentes colores (RB o BR):
$\frac25 \times \frac34 + \frac35 \times \frac24
= \frac{6}{20} + \frac{6}{20} = \frac{12}{20} = \frac35$.
\end{solution}

\begin{solution}{exo:g9:statproba:8}
\emph{1.} $9 \times 14 = 126$ agujas.

\emph{2.} un promedio de $15$ encima $10$ coincidencias significa $150$ puntos en
total: ella debe anotar $150 - 126 = 24$ agujas.
\end{solution}

\begin{solution}{exo:g9:statproba:9}
\emph{1.} entre los $36$ pares igualmente probables, los dobles son
$(1,1), (2,2), \dots, (6,6)$: seis de ellos, entonces
$\P(\text{double}) = \frac{6}{36} = \frac16$.

\emph{2.} Dos juegos independientes: árbol con victoria ($\frac16$) / perder
($\frac56$) en cada nivel. \omterm{def:g9:statproba:probability}{Probabilidad} de nunca ganar:
$\frac56 \times \frac56 = \frac{25}{36}$. Entonces
\[
\P(\text{at least one win}) = 1 - \frac{25}{36} = \frac{11}{36}
\approx 0.31 .
\]
\end{solution}

\begin{solution}{pb:g9:statproba:1}
\textbf{1.} Cada rollo tiene $5$ no seis resultados de $6$; el
$6 \times 6$ la tabla muestra $5 \times 5 = 25$ seis resultados libres
de $36$: $ P(\text{no six}) = \frac{25}{36}$, entonces
$ P(\text{at least one}) = 1 - \frac{25}{36} = \frac{11}{36}
\approx 0.31$.

\textbf{2.} $\frac13 = \frac{12}{36} \neq \frac{11}{36}$. el
La suma cuenta el resultado "seis y luego seis" una vez en cada
$\frac16$: una vez con demasiada frecuencia. Probabilidades de eventos que pueden
suceder \emph{juntos} no simplemente agregue.

\textbf{3.} \omterm{def:g1:addition:def}{Suma} $7$: los seis resultados $(1,6), (2,5), (3,4),
(4,3), (5,2), (6,1)$: $ P = \frac{6}{36} = \frac16$. \omterm{def:g1:addition:def}{Suma} $2$: solo $(1,1)$: $\frac{1}{36}$. Siete tiene la mayor cantidad de formas de
suceder --- la mitad gorda de la \omterm{def:g1:addition:def}{sumas} mesa.

\textbf{4.} \omterm{def:g1:addition:def}{Suma} $9$: $(3,6), (4,5), (5,4), (6,3)$ --- cuatro
resultados, $\frac{4}{36}$. \omterm{def:g1:addition:def}{Suma} $10$: $(4,6), (5,5), (6,4)$ ---
tres, $\frac{3}{36}$. Nueve es más probable: el conde decide,
no el número de ``maneras de escribir'' el \omterm{def:g1:addition:def}{suma} con desordenado
\omterm{def:g5:solids:def}{caras}.

\textbf{5.} Cada nivel del árbol multiplica el número seis.
\omterm{def:g9:statproba:probability}{probabilidad} por $\frac56$, lo que pasó antes: después $ n $
rollos, $ P(\text{no six at all}) = \left(\frac56\right)^n $, y
la regla del complemento da la fórmula.

\textbf{6.} $\left(\frac56\right)^4 = \frac{625}{1296} \approx
0.482$, por lo que el Chevalier ganó con \omterm{def:g9:statproba:probability}{probabilidad}
$1 - 0.482 = 0.518$ --- un borde silencioso y constante de casi
$2\,\%$ por partido: una excelente apuesta, y pagó por su
carruajes.

\textbf{7.} En $7$ tira su regla da $\frac76$ --- a
\omterm{def:g9:statproba:probability}{probabilidad} más que $1$, tonterías. es la pregunta 2
conteo doble, compuesto: los éxitos de las tiradas se superponen y
sumando sus posibilidades cuenta las superposiciones una y otra vez.

\textbf{8.} $1 - \left(\frac{35}{36}\right)^{24} \approx
1 - 0.5086 = 0.491$: debajo de uno \omterm{def:g3:division:half}{medio}. Apuesta \omterm{def:g3:numbers:evenodd}{incluso} dinero en un
$49.1\,\%$ evento, el Chevalier perdió aproximadamente $2$ juegos en cada
$100$ en promedio --- lenta, misteriosamente (para él), y
inevitablemente.

\textbf{9.} $1 - \left(\frac{35}{36}\right)^{25} \approx
1 - 0.494 = 0.506$: con $25$ tira la apuesta se vuelve favorable.
Una tirada separó al Chevalier de la rentabilidad.

\textbf{10.} multiplicando un \omterm{def:g9:statproba:probability}{probabilidad} por el número de intentos
cuenta los éxitos superpuestos varias veces y, finalmente,
produce respuestas imposibles arriba $1$. La ruta correcta es
siempre a través del complemento: encadenar las probabilidades de fallo
por multiplicación, luego restar de $1$.

\textbf{11.} \omterm{def:g9:statproba:indicators}{Significar}: $\frac{9 \times 2\,000 + 20\,000}{10} =
\frac{38\,000}{10} = 3\,800$ euros. \omterm{def:g9:statproba:indicators}{Mediana}: los salarios medios
son ambos $2\,000$: \omterm{def:g9:statproba:indicators}{mediana} $2\,000$ euros. el anuncio
debería citar el \omterm{def:g9:statproba:indicators}{mediana} --- nueve de los diez trabajadores nunca ven
algo parecido $3\,800$; el \omterm{def:g9:statproba:indicators}{significar} es arrastrado por uno
valor atípico (\cref{ex:g9:statproba:weighted} advertido de pesas).

\textbf{12.} Por ejemplo $2, 8, 15, 17, 18$: \omterm{def:g9:statproba:indicators}{mediana} $15$
(valor medio), \omterm{def:g9:statproba:indicators}{significar} $\frac{60}{5} = 12$. A \omterm{def:g9:statproba:indicators}{significar} debajo del
\omterm{def:g9:statproba:indicators}{mediana} traiciona una cola de \emph{bajo} marcas tirando del promedio
abajo mientras que la parte superior \omterm{def:g3:division:half}{medio} se sienta alto.

\textbf{13.} Primeros tres factores:
$1 \times \frac{364}{365} \times \frac{363}{365} \approx
0.992$. Cada nuevo estudiante debe esquivar una fecha más tomada, por lo que el
factores se reducen: $\frac{343}{365} \approx 0.94$ para el
$23$ rd. el completo \omterm{def:g2:mult:def}{producto} cae a aproximadamente $0.493$ --- entonces
\[
P(\text{at least one shared birthday}) \approx 0.507 :
\]
mejor que \omterm{def:g3:numbers:evenodd}{incluso}, en una clase de sólo $23$. La intuición dice
``$23$ gente, $365$ días, no hay posibilidad''; las matemáticas dicen
"Lanzar una moneda".

\textbf{14.} $\frac{23 \times 22}{2} = 253$ pares. Cada par es
una nueva oportunidad para una coincidencia, y $253$ oportunidades
en $\frac{1}{365}$ cada uno ya no es un asunto menor: el
la sorpresa se disuelve una vez que se cuentan las parejas, no las personas.

\textbf{15.} Protocolo de ejemplo: tirar un dado en bloques de cuatro,
registrar si cada bloque muestra un seis, por ejemplo $100$ bloques; o
recopilar las listas de cumpleaños de muchas clases de aproximadamente $23$ y
registre la proporción con una coincidencia. lo observado
\omterm{def:g7:stats:frequency}{frecuencia} se tambaleará, pero a medida que aumente el número de repeticiones
debería establecerse cada vez más cerca de las probabilidades calculadas
($0.518$; $0.507$) --- las frecuencias convergen en probabilidades:
la ley de los grandes números, expresada aquí como una expectativa y
demostrado en el volumen de High School.
\end{solution}
